\begin{center}
    \Large{\textbf{Аннотация}}
\end{center}

Исследуется сходимость градиентных методов оптимизации первого порядка с предобуславливанием, использующих регуляризацию с затуханием весов.
Особое внимание уделяется популярным методам из данного класса, таким как AdamW и OASIS.
Рассматриваются три способа добавления регуляризации в методы с предобуславливанием, включая новый, промежуточный способ (масштабированное затухание весов), и изучается вопрос о том, к решению какой задачи сходится каждый из них.
Показывается, что методы с затуханием весов оптимизируют не исходную регуляризованную функцию потерь, а изменённую функцию с адаптивным покоординатным регуляризатором, и доказывается нижняя оценка расстояния между решениями этих двух задач.
Доказываются четыре теоремы о скорости сходимости методов с предобуславливанием и затуханием весов: в невыпуклом и сильно выпуклом случаях, каждый --- в детерминированной и стохастической постановках.
Проводятся вычислительные эксперименты на эталонных наборах данных и моделях, включая количественную проверку теоретических предсказаний и анализ гиперпараметров, чтобы сравнить рассматриваемые методы на реальных задачах.
